\chapter{Predizioni Rischiose Datate: Falsificabilità
Pubblica}\label{predizioni-rischiose-datate-falsificabilituxe0-pubblica}

\begin{quote}
\textbf{Argomento portante:} una teoria sociale che non si espone alla
falsificazione \emph{datata} è indistinguibile da una narrazione. Questo
capitolo aggrega le sedici predizioni rischiose della DEQ$^2$ --- tredici
derivate dai Cap. 7-9 sui tre attori sistemici (cinque USA, quattro
Cina, quattro Brasile), tre derivate dal Cap. 10 sull'aggregato globale
--- in un \emph{pre-registered prediction set} con soglie numeriche,
date di scadenza, fonti dati pubbliche e protocollo di valutazione. Per
ciascuna predizione è inoltre specificato \emph{cosa accadrebbe alla
teoria se fallisse}: la falsificazione di alcune predizioni è gestibile
come correzione locale, quella di altre forzerebbe la riformulazione
strutturale dell'apparato.
\end{quote}



\section{Criteri DEQ$^2$ per una Predizione
Rischiosa}\label{criteri-dequxb2-per-una-predizione-rischiosa}

\Hypoth{} \textbf{Eredità popperiana operativizzata.} Una predizione
\(\mathcal{P}\) è \emph{rischiosa} nel senso DEQ$^2$ se soddisfa
simultaneamente:

\begin{enumerate}
\def\labelenumi{\arabic{enumi}.}
\tightlist
\item
  \textbf{Datazione esplicita}: esiste \(t^*\) entro cui \(\mathcal{P}\)
  è valutabile;
\item
  \textbf{Soglia numerica}: esiste un valore di confronto \(v^*\) tale
  che \(\mathcal{P}\) asserisce \(X(t^*) \lessgtr v^*\);
\item
  \textbf{Fonte dati pubblica}: esiste una sorgente \(\mathcal{D}\)
  verificabile da terzi indipendenti;
\item
  \textbf{Probabilità soggettiva esplicita}:
  \(\Pr(\mathcal{P} \text{ vera}) = p\) è dichiarata, \(p \in (0, 1)\);
\item
  \textbf{Conseguenza teorica}: è dichiarato cosa accade alla DEQ$^2$ se
  \(\mathcal{P}\) fallisce --- \emph{correzione locale},
  \emph{aggiustamento parametrico}, oppure \emph{riformulazione
  strutturale}.
\end{enumerate}

\Proved{} \textbf{Risultato metodologico.} Una predizione che soddisfa
(1)-(5) è \emph{Bayes-aggiornabile}: la sua valutazione produce un
aggiornamento esplicito delle probabilità posteriori sui parametri della
teoria. La DEQ$^2$ si presenta come \emph{modello probabilistico vivente},
non come dottrina chiusa.



\section{Le Tredici Predizioni dei Cap.
7-9}\label{le-tredici-predizioni-dei-cap.-7-9}

\subsection{USA --- Egemone Decoerente (Cap.
7)}\label{usa-egemone-decoerente-cap.-7}

{\def\LTcaptype{none} % do not increment counter
{\small\begin{longtable}[]{@{}
  >{\raggedright\arraybackslash}p{(\linewidth - 12\tabcolsep) * \real{0.1429}}
  >{\raggedright\arraybackslash}p{(\linewidth - 12\tabcolsep) * \real{0.1429}}
  >{\raggedright\arraybackslash}p{(\linewidth - 12\tabcolsep) * \real{0.1429}}
  >{\raggedright\arraybackslash}p{(\linewidth - 12\tabcolsep) * \real{0.1429}}
  >{\raggedright\arraybackslash}p{(\linewidth - 12\tabcolsep) * \real{0.1429}}
  >{\raggedright\arraybackslash}p{(\linewidth - 12\tabcolsep) * \real{0.1429}}
  >{\raggedright\arraybackslash}p{(\linewidth - 12\tabcolsep) * \real{0.1429}}@{}}
\toprule\noalign{}
\begin{minipage}[b]{\linewidth}\raggedright
\#
\end{minipage} & \begin{minipage}[b]{\linewidth}\raggedright
Predizione
\end{minipage} & \begin{minipage}[b]{\linewidth}\raggedright
\(t^*\)
\end{minipage} & \begin{minipage}[b]{\linewidth}\raggedright
\(v^*\)
\end{minipage} & \begin{minipage}[b]{\linewidth}\raggedright
Fonte \(\mathcal{D}\)
\end{minipage} & \begin{minipage}[b]{\linewidth}\raggedright
\(p\)
\end{minipage} & \begin{minipage}[b]{\linewidth}\raggedright
Falsificazione
\end{minipage} \\
\midrule\noalign{}
\endhead
\bottomrule\noalign{}
\endlastfoot
\textbf{U1} & V-Dem Liberal Democracy Index USA \(< 0{,}55\) & report
aprile 2027 (anno 2026) & \(0{,}55\) & \url{https://v-dem.net} &
\(0{,}70\) & locale (parametro \(\Phi\)) \\
\textbf{U2} & Quota dollaro nelle riserve IMF COFER \(< 55\%\) & Q4 2027
& \(0{,}55\) & IMF COFER trimestrale & \(0{,}55\) & locale (parametro
\(\sigma\) esogena) \\
\textbf{U3} & Quota dollaro \(< 52\%\) & Q4 2030 & \(0{,}52\) & IMF
COFER trimestrale & \(0{,}50\) & locale (cumulativa con U2) \\
\textbf{U4} & Vetoes presidenziali \(> 15\) + emergency docket SCOTUS
\(> 30\) in 12 mesi & gennaio-luglio 2027, primo anno Congresso 119$^\circ$ &
\(15\) vetoes, \(30\) emergency & Federal Register + SCOTUSblog &
\(0{,}55\) & parametrica (\(\Omega\) acuta vs cronica) \\
\textbf{U5} & Correlazione Musk visibility--polarizzazione USA
\(r > 0{,}3\) & calibrazione 2022-2026 & \(r = 0{,}3\) & dataset
compilato (Cap. 7) & \(0{,}65\) & locale (modello
\(\zeta_{\text{Sol}}\)) \\
\end{longtable}}
}

\subsection{Cina --- Rigido Sincronizzato (Cap.
8)}\label{cina-rigido-sincronizzato-cap.-8}

{\def\LTcaptype{none} % do not increment counter
{\small\begin{longtable}[]{@{}
  >{\raggedright\arraybackslash}p{(\linewidth - 12\tabcolsep) * \real{0.1429}}
  >{\raggedright\arraybackslash}p{(\linewidth - 12\tabcolsep) * \real{0.1429}}
  >{\raggedright\arraybackslash}p{(\linewidth - 12\tabcolsep) * \real{0.1429}}
  >{\raggedright\arraybackslash}p{(\linewidth - 12\tabcolsep) * \real{0.1429}}
  >{\raggedright\arraybackslash}p{(\linewidth - 12\tabcolsep) * \real{0.1429}}
  >{\raggedright\arraybackslash}p{(\linewidth - 12\tabcolsep) * \real{0.1429}}
  >{\raggedright\arraybackslash}p{(\linewidth - 12\tabcolsep) * \real{0.1429}}@{}}
\toprule\noalign{}
\begin{minipage}[b]{\linewidth}\raggedright
\#
\end{minipage} & \begin{minipage}[b]{\linewidth}\raggedright
Predizione
\end{minipage} & \begin{minipage}[b]{\linewidth}\raggedright
\(t^*\)
\end{minipage} & \begin{minipage}[b]{\linewidth}\raggedright
\(v^*\)
\end{minipage} & \begin{minipage}[b]{\linewidth}\raggedright
Fonte \(\mathcal{D}\)
\end{minipage} & \begin{minipage}[b]{\linewidth}\raggedright
\(p\)
\end{minipage} & \begin{minipage}[b]{\linewidth}\raggedright
Falsificazione
\end{minipage} \\
\midrule\noalign{}
\endhead
\bottomrule\noalign{}
\endlastfoot
\textbf{C1} & Deflazione cinese \(> 12\) trimestri consecutivi & Q4 2026
& \(12\) trim. & NBS / PIIE / WB & \(0{,}65\) & locale (\(\Omega\)
misurato) \\
\textbf{C2} & \(P(\text{evento *bloqueo* Taiwan})\) in 2027-2029,
condizionale a (deflazione \(\geq 12\) trim.) \(\wedge\) (GDP \(< 4\%\))
\(\wedge\) (nascite \(< 5/1000\)) & 2027-2029 &
\(P \in [0{,}25, 0{,}40]\) & corso eventi + INSS / AEI / Crisis Group &
\(0{,}50\) entro intervallo & strutturale se \(P\) fuori intervallo \\
\textbf{C3} & Yuan in IMF COFER \(< 4\%\) nel 2027 e \(< 5\%\) nel 2030
& Q4 2027, Q4 2030 & \(4\%\), \(5\%\) & IMF COFER trimestrale &
\(0{,}60\) & locale (modello dedollarizzazione) \\
\textbf{C4} & PPI cinese in territorio negativo \(\geq 5\) anni cumulati
& Q4 2027 & \(5\) anni & NBS PPI mensile & \(0{,}70\) & locale
(anti-involution efficacia) \\
\end{longtable}}
}

\subsection{Brasile --- Bilanciere Quantistico (Cap.
9)}\label{brasile-bilanciere-quantistico-cap.-9}

{\def\LTcaptype{none} % do not increment counter
{\small\begin{longtable}[]{@{}
  >{\raggedright\arraybackslash}p{(\linewidth - 12\tabcolsep) * \real{0.1429}}
  >{\raggedright\arraybackslash}p{(\linewidth - 12\tabcolsep) * \real{0.1429}}
  >{\raggedright\arraybackslash}p{(\linewidth - 12\tabcolsep) * \real{0.1429}}
  >{\raggedright\arraybackslash}p{(\linewidth - 12\tabcolsep) * \real{0.1429}}
  >{\raggedright\arraybackslash}p{(\linewidth - 12\tabcolsep) * \real{0.1429}}
  >{\raggedright\arraybackslash}p{(\linewidth - 12\tabcolsep) * \real{0.1429}}
  >{\raggedright\arraybackslash}p{(\linewidth - 12\tabcolsep) * \real{0.1429}}@{}}
\toprule\noalign{}
\begin{minipage}[b]{\linewidth}\raggedright
\#
\end{minipage} & \begin{minipage}[b]{\linewidth}\raggedright
Predizione
\end{minipage} & \begin{minipage}[b]{\linewidth}\raggedright
\(t^*\)
\end{minipage} & \begin{minipage}[b]{\linewidth}\raggedright
\(v^*\)
\end{minipage} & \begin{minipage}[b]{\linewidth}\raggedright
Fonte \(\mathcal{D}\)
\end{minipage} & \begin{minipage}[b]{\linewidth}\raggedright
\(p\)
\end{minipage} & \begin{minipage}[b]{\linewidth}\raggedright
Falsificazione
\end{minipage} \\
\midrule\noalign{}
\endhead
\bottomrule\noalign{}
\endlastfoot
\textbf{B1} & \(P(\text{Lula o successore di sinistra vince elezioni})\)
aprile 2026 & aprile, giugno, settembre 2026 (revisione mobile) &
\(p \in [0{,}55, 0{,}70]\) & Genial/Quaest, AtlasIntel, Datafolha &
meta-stima soggettiva & aggiustamento parametrico al cambiare
evidenze \\
\textbf{B2} & Se A/B vince: TFFF impegni \(> 5\)B USD entro fine 2027 e
BRICS BMG operativo entro Q2 2027 & Q4 2027, Q2 2027 & \(5\)B,
\emph{operativo} & comunicati BRICS, World Bank & \(0{,}55\)
condizionale & locale (operatività iniziativa) \\
\textbf{B3} & Se A/B vince: V-Dem Brasile Liberal Democracy Index
\(> 0{,}65\) & report 2028 (anno 2027) & \(0{,}65\) &
\url{https://v-dem.net} & \(0{,}65\) condizionale & locale (\(\Phi\)
Brasile) \\
\textbf{B4} & Se C vince: deterioramento V-Dem \(> 0{,}05\) in 24 mesi +
uscita da Group of Friends for Peace + BRICS partecipazione passiva & 24
mesi post-insediamento (gen 2027 $\to$ gen 2029) & \(\Delta \geq 0{,}05\)
punti, evento ufficiale, evento osservabile & V-Dem + comunicati
ufficiali & \(0{,}60\) condizionale & strutturale (la
\emph{solitonizzazione} discorsiva è il meccanismo principale del Cap.
7.3) \\
\end{longtable}}
}

\Proved{} \textbf{Sintesi delle 13 predizioni Cap. 7-9.} Probabilità
soggettiva mediana: \(p \approx 0{,}60\). Distribuzione: 2 strutturali
(C2, B4), 2 di meta-stima (B1, U5), 9 locali. Il fallimento di una
predizione locale aggiusterebbe parametri specifici; il fallimento di
una predizione strutturale costringerebbe a ripensare l'apparato
corrispondente.



\section{Tre Predizioni di Sistema Globale (Cap.
10)}\label{tre-predizioni-di-sistema-globale-cap.-10}

{\def\LTcaptype{none} % do not increment counter
{\small\begin{longtable}[]{@{}
  >{\raggedright\arraybackslash}p{(\linewidth - 12\tabcolsep) * \real{0.1429}}
  >{\raggedright\arraybackslash}p{(\linewidth - 12\tabcolsep) * \real{0.1429}}
  >{\raggedright\arraybackslash}p{(\linewidth - 12\tabcolsep) * \real{0.1429}}
  >{\raggedright\arraybackslash}p{(\linewidth - 12\tabcolsep) * \real{0.1429}}
  >{\raggedright\arraybackslash}p{(\linewidth - 12\tabcolsep) * \real{0.1429}}
  >{\raggedright\arraybackslash}p{(\linewidth - 12\tabcolsep) * \real{0.1429}}
  >{\raggedright\arraybackslash}p{(\linewidth - 12\tabcolsep) * \real{0.1429}}@{}}
\toprule\noalign{}
\begin{minipage}[b]{\linewidth}\raggedright
\#
\end{minipage} & \begin{minipage}[b]{\linewidth}\raggedright
Predizione
\end{minipage} & \begin{minipage}[b]{\linewidth}\raggedright
\(t^*\)
\end{minipage} & \begin{minipage}[b]{\linewidth}\raggedright
\(v^*\)
\end{minipage} & \begin{minipage}[b]{\linewidth}\raggedright
Fonte \(\mathcal{D}\)
\end{minipage} & \begin{minipage}[b]{\linewidth}\raggedright
\(p\)
\end{minipage} & \begin{minipage}[b]{\linewidth}\raggedright
Falsificazione
\end{minipage} \\
\midrule\noalign{}
\endhead
\bottomrule\noalign{}
\endlastfoot
\textbf{G1} &
\(P(\text{sistema globale 2031 in regime di decoerenza generalizzata})\)
--- operazionalizzato come \emph{almeno due fra USA, Cina, UE in
\(\mathcal{R}_-\)} & dicembre 2031 & \(\geq 2\) attori in
\(\mathcal{R}_-\) & aggregazione V-Dem, Pew, Brand Finance, IMF, NBS al
2031 & \(0{,}80\) & strutturale (forma di \(T_s^{(2),\text{sys}}\)) \\
\textbf{G2} &
\(P(\text{sistema globale 2031 con almeno un $\Phi$-engine globale operativo})\)
--- operazionalizzato come Brasile \(C_{\text{DEQ}}^2 \geq 1\)
\emph{oppure} nuovo attore equivalente & dicembre 2031 &
\(C_{\text{DEQ}}^2 \geq 1\) per almeno un attore & calcolo dashboard
SPC-DEQ & \([0{,}55, 0{,}70]\), condizionale a B1 & locale (dipende da
B1) \\
\textbf{G3} & Cina: divergenza \(d^{\text{oss}} - d^{\text{spont}}\)
crescente, con \emph{fatica del \(K\) imposto} visibile & finestra di
osservazione 2029-2031 &
\(\Delta(d^{\text{oss}} - d^{\text{spont}}) > 0{,}10\) in 24 mesi &
dashboard SPC-DEQ + V-Dem + NBS & \(0{,}55\) & strutturale (definizione
di \emph{Rigido Sincronizzato} del Cap. 8.0) \\
\end{longtable}}
}

\Hypoth{} \textbf{Lettura.} G1 e G2 sono \emph{complementari}: la teoria
afferma simultaneamente che il sistema globale 2031 sarà in decoerenza
\emph{e} che esiste con probabilità \(\approx 0{,}55\)-\(0{,}70\) un
\(\Phi\)-engine operativo. La conferma di entrambe
(\(\Pr \approx 0{,}50\)) consoliderebbe il quadro M1 del Cap. 10.4. La
conferma di G1 con falsificazione di G2 (\(\Pr \approx 0{,}25\))
realizzerebbe M2 --- multipolarità degenere senza centro coerente.



\section{La Predizione Meta: l'Elezione Brasiliana del 4 Ottobre
2026}\label{la-predizione-meta-lelezione-brasiliana-del-4-ottobre-2026}

\Hypoth{} \textbf{Tesi.} B1 non è una predizione \emph{fra le altre} ---
è la \textbf{predizione che condiziona tutte quelle brasiliane (B2, B3,
B4) e una predizione globale (G2)}. La sua valutazione, prevista per il
4 ottobre 2026 (primo turno) e l'ottobre/novembre 2026 (eventuale
secondo turno), è il \emph{singolo evento più informativo} dell'intero
apparato DEQ$^2$ fra aprile 2026 e dicembre 2031.

\Proved{} \textbf{Protocollo di valutazione di B1.}

{\def\LTcaptype{none} % do not increment counter
{\small\begin{longtable}[]{@{}
  >{\raggedright\arraybackslash}p{(\linewidth - 4\tabcolsep) * \real{0.3333}}
  >{\raggedright\arraybackslash}p{(\linewidth - 4\tabcolsep) * \real{0.3333}}
  >{\raggedright\arraybackslash}p{(\linewidth - 4\tabcolsep) * \real{0.3333}}@{}}
\toprule\noalign{}
\begin{minipage}[b]{\linewidth}\raggedright
Data
\end{minipage} & \begin{minipage}[b]{\linewidth}\raggedright
Azione
\end{minipage} & \begin{minipage}[b]{\linewidth}\raggedright
Conseguenza
\end{minipage} \\
\midrule\noalign{}
\endhead
\bottomrule\noalign{}
\endlastfoot
30 giugno 2026 & Re-stima \(p_{B1}\) su evidenze sondaggi aggregati &
aggiornamento bayesiano della probabilità del Mondo M1/M2 \\
30 settembre 2026 & Re-stima \(p_{B1}\) alla vigilia del primo turno &
aggiornamento finale pre-evento \\
4 ottobre 2026, sera & Risultato primo turno: \(\geq 50\%\) Lula? Se sì
$\to$ B1 vera & falsificazione binaria \\
Ottobre/novembre 2026 (eventuale 2$^\circ$ turno) & Risultato secondo turno:
vince A/B o C? & falsificazione binaria definitiva \\
Entro 30 novembre 2026 & Pubblicazione del \emph{post-mortem} B1 &
aggiornamento di tutta la cascata B2-B3-B4-G2 \\
\end{longtable}}
}

\Hypoth{} \textbf{Conseguenza per la teoria.} Se B1 = \emph{vera}
(vittoria A/B): la cascata B2, B3, G2 conserva la sua probabilità
condizionale. Se B1 = \emph{falsa} (vittoria C): la cascata B2, B3, G2
si riorienta verso B4 e M2. In entrambi i casi, la DEQ$^2$ resta valida ---
ma il \emph{contenuto empirico futuro} del libro cambia
significativamente.



\section{Protocollo di Pubblicazione e
Falsificazione}\label{protocollo-di-pubblicazione-e-falsificazione}

\Hypoth{} \textbf{Pre-registrazione.} Le sedici predizioni di questo
capitolo sono \emph{pre-registrate} alla data di pubblicazione del libro
(target 2026 Q3 --- KDP IT/EN/ES/PT). La pre-registrazione consiste
nella pubblicazione contestuale di:

\begin{enumerate}
\def\labelenumi{\arabic{enumi}.}
\tightlist
\item
  \textbf{Tabella delle predizioni} (sezioni 11.1, 11.2) con tutti i
  campi (1)-(5);
\item
  \textbf{Repository pubblico} (GitHub
  \texttt{bilanciere-quantistico/predizioni-deq2}) contenente:

  \begin{itemize}
  \tightlist
  \item
    file CSV con tutte le predizioni e i loro stati (pendente / vera /
    falsa);
  \item
    script Python di calcolo automatico delle predizioni quantitative
    dai dataset pubblici;
  \item
    log dei \emph{checkpoint} trimestrali con re-stima probabilità
    soggettive;
  \end{itemize}
\item
  \textbf{Protocollo di aggiornamento}: ogni 31 marzo dell'anno
  successivo, pubblicazione di un \emph{update report} con:

  \begin{itemize}
  \tightlist
  \item
    stato di ciascuna predizione;
  \item
    aggiornamento bayesiano delle probabilità;
  \item
    eventuali correzioni alla teoria (locali/parametriche/strutturali).
  \end{itemize}
\end{enumerate}

\Proved{} \textbf{Conseguenza editoriale.} Il libro non è
\emph{autorevole per autorevolezza} ma \emph{autorevole per esposizione
al rischio}. Ogni copia venduta riceve l'URL del repository e ha accesso
al ciclo di update. Questo è coerente con la posizione epistemologica
della Topografia App. A: produrre segnali \((x, y, z)\) alti su
\emph{tutti} gli assi --- incluso \(z\) (complessità dialettica) tramite
l'esposizione esplicita alla falsificazione.



\section{Cosa Accadrebbe se le Predizioni
Fallissero}\label{cosa-accadrebbe-se-le-predizioni-fallissero}

\Hypoth{} \textbf{Mappa di vulnerabilità della teoria.}

{\def\LTcaptype{none} % do not increment counter
{\small\begin{longtable}[]{@{}
  >{\raggedright\arraybackslash}p{(\linewidth - 2\tabcolsep) * \real{0.5000}}
  >{\raggedright\arraybackslash}p{(\linewidth - 2\tabcolsep) * \real{0.5000}}@{}}
\toprule\noalign{}
\begin{minipage}[b]{\linewidth}\raggedright
Predizione
\end{minipage} & \begin{minipage}[b]{\linewidth}\raggedright
Fallimento $\to$ Conseguenza per DEQ$^2$
\end{minipage} \\
\midrule\noalign{}
\endhead
\bottomrule\noalign{}
\endlastfoot
U1 & V-Dem USA migliora rapidamente: aggiusta \(\Phi^{\text{USA}}\) e
\(\zeta_{\text{Sol}}\) effettivo. Locale. \\
U2-U3 & Dollaro tiene: rivedere stima \(\sigma\) esogena globale.
Locale. \\
U4 & Bassa attività veto/SCOTUS: rivedere stima \(\Omega^{\text{USA}}\).
Locale. \\
U5 & Correlazione \(r < 0{,}3\): il modello \(\zeta_{\text{Sol}}(t)\)
del Cap. 6.6 va riformulato. Locale-parametrica. \\
C1 & Deflazione finisce prima: anti-involution funziona. Aggiustamento
\(\Omega^{\text{Cina}}\). Locale. \\
C2 & Nessun bloqueo malgrado condizioni soddisfatte: ipotesi
\emph{Taiwan come valvola} va rivista. \textbf{Strutturale}. \\
C3 & Yuan oltre \(5\%\) entro 2030: la dedollarizzazione converge sullo
yuan, contro la tesi Cap. 8.6. Locale-parametrica. \\
C4 & PPI positivo prima del 2027: anti-involution efficace.
Aggiustamento \(\Omega^{\text{Cina}}\). Locale. \\
B1 & Vittoria di Flávio o destra radicale: B2-B3 cadono, attiva ramo B4.
Cascata controfattuale prevista. \\
B2-B3 & (condizionali a B1 vera) fallimento: rivedere capacità operativa
BRICS+ e \(s_{iii}^{\text{Brasile}}\). Locale. \\
B4 & (condizionale a B1 falsa) fallimento --- Brasile resiste sotto
destra radicale: la \emph{solitonizzazione} del Cap. 7.3 va ripensata
come meccanismo non automatico. \textbf{Strutturale}. \\
G1 & Solo uno o nessun grande attore in \(\mathcal{R}_-\) entro 2031: la
forma stessa di \(T_s^{(2),\text{sys}}\) va ricalibrata.
\textbf{Strutturale}. \\
G2 & Brasile fallisce \(C_{\text{DEQ}}^2 \geq 1\) malgrado A/B: la
definizione di \(\Phi\)-engine è insufficiente. \textbf{Strutturale}. \\
G3 & Cina mantiene \(d^{\text{oss}} \approx d^{\text{spont}}\): la
categoria \emph{Rigido Sincronizzato} non descrive la Cina come
ipotizzato. \textbf{Strutturale}. \\
\end{longtable}}
}

\Proved{} \textbf{Sintesi.} Su 16 predizioni: \textbf{5 strutturali}
(C2, B4, G1, G2, G3), \textbf{9 locali}, \textbf{2 di meta-stima} (B1,
U5). Il fallimento di una predizione strutturale impone
\emph{riformulazione}; il fallimento di tutte le strutturali
simultaneamente \emph{invaliderebbe la teoria nel suo insieme}.
Probabilità di fallimento simultaneo (assumendo indipendenza
condizionale debole): \(\Pr \approx (1-0{,}55)^5 \approx 0{,}018\) ---
basso ma non trascurabile.



\section{Sintesi del Capitolo}\label{sintesi-del-capitolo}

\Proved{} \textbf{In una frase:} la DEQ$^2$ si espone deliberatamente a
sedici predizioni datate, con probabilità soggettive esplicite e fonti
dati pubbliche; cinque di queste sono strutturali e potrebbero
invalidare l'apparato; il libro è dunque un \emph{pre-registered
prediction set} il cui valore epistemico cresce o diminuisce in
proporzione alla performance verificabile della teoria.

{\def\LTcaptype{none} % do not increment counter
{\small\begin{longtable}[]{@{}
  >{\raggedright\arraybackslash}p{(\linewidth - 6\tabcolsep) * \real{0.2500}}
  >{\raggedright\arraybackslash}p{(\linewidth - 6\tabcolsep) * \real{0.2500}}
  >{\raggedright\arraybackslash}p{(\linewidth - 6\tabcolsep) * \real{0.2500}}
  >{\raggedright\arraybackslash}p{(\linewidth - 6\tabcolsep) * \real{0.2500}}@{}}
\toprule\noalign{}
\begin{minipage}[b]{\linewidth}\raggedright
Categoria
\end{minipage} & \begin{minipage}[b]{\linewidth}\raggedright
Numero
\end{minipage} & \begin{minipage}[b]{\linewidth}\raggedright
Probabilità mediana \(p\)
\end{minipage} & \begin{minipage}[b]{\linewidth}\raggedright
Tipo di falsificazione
\end{minipage} \\
\midrule\noalign{}
\endhead
\bottomrule\noalign{}
\endlastfoot
Cap. 7 USA & 5 & \(0{,}60\) & 4 locali, 1 parametrica/meta (U5) \\
Cap. 8 Cina & 4 & \(0{,}60\) & 3 locali, 1 strutturale (C2) \\
Cap. 9 Brasile & 4 & meta + condizionali & 2 condizionali, 1 strutturale
(B4), 1 meta (B1) \\
Cap. 10 globali & 3 & \(0{,}65\) & 2 strutturali (G1, G2 e G3 entrambe),
1 condizionale \\
\textbf{Totale} & \textbf{16} & \(\approx 0{,}60\) & \textbf{5
strutturali, 9 locali/condizionali, 2 meta} \\
\end{longtable}}
}

\Hypoth{} \textbf{Tesi conclusiva.} Le predizioni di questo capitolo
sono il \emph{contenuto empirico futuro} della DEQ$^2$. La Parte V del
libro (Cap. 12-13) costruirà invece il \emph{contenuto normativo}: dato
lo spazio degli scenari probabilizzati, quali sono le strategie
antifragili --- individuali, organizzative, statali --- coerenti con
ciascuno dei quattro mondi possibili al 2031?



\section{Note Epistemiche di Chiusura
Capitolo}\label{note-epistemiche-di-chiusura-capitolo}

{\def\LTcaptype{none} % do not increment counter
{\small\begin{longtable}[]{@{}
  >{\raggedright\arraybackslash}p{(\linewidth - 4\tabcolsep) * \real{0.3333}}
  >{\raggedright\arraybackslash}p{(\linewidth - 4\tabcolsep) * \real{0.3333}}
  >{\raggedright\arraybackslash}p{(\linewidth - 4\tabcolsep) * \real{0.3333}}@{}}
\toprule\noalign{}
\begin{minipage}[b]{\linewidth}\raggedright
\#
\end{minipage} & \begin{minipage}[b]{\linewidth}\raggedright
Affermazione
\end{minipage} & \begin{minipage}[b]{\linewidth}\raggedright
Status
\end{minipage} \\
\midrule\noalign{}
\endhead
\bottomrule\noalign{}
\endlastfoot
1 & Cinque criteri DEQ$^2$ per predizione rischiosa & \Hypoth{}
metodologia \\
2 & Tredici predizioni Cap. 7-9 aggregate con campi (1)-(5) & \Proved{}
aggregazione completa \\
3 & Tre predizioni di sistema globale dal Cap. 10 & \Hypoth{}
derivazione coerente \\
4 & B1 come predizione meta che condiziona la cascata & \Hypoth{}
architettura logica \\
5 & Protocollo di pre-registrazione e repository GitHub & \Hypoth{}
commitment editoriale \\
6 & Mappa di vulnerabilità delle 16 predizioni & \Hypoth{} analisi
strutturale \\
7 & \(\Pr(\text{fallimento simultaneo strutturale}) \approx 0{,}018\) &
\Conj{} stima sotto indipendenza condizionale \\
8 & \(5\) strutturali + \(9\) locali/condizionali + \(2\) meta &
\Proved{} classificazione completa \\
9 & Update report annuale al 31 marzo & \Hypoth{} protocollo
operativo \\
\end{longtable}}
}



\section{Chiusura della Parte IV e Apertura della Parte
V}\label{chiusura-della-parte-iv-e-apertura-della-parte-v}

La \textbf{Parte IV} è ora completa. Cap. 10 ha tradotto l'apparato in
dashboard operativa (sei pannelli + aggregato + cinque regole WE-DEQ$^2$ +
indice \(C_{\text{DEQ}}^2\) + scenari + quattro mondi possibili 2031).
Cap. 11 ha aggregato le 16 predizioni rischiose datate con criteri
popperiani operativi e protocollo di falsificazione pubblica.

La \textbf{Parte V} può ora costruire il contenuto normativo:

\begin{itemize}
\tightlist
\item
  \textbf{Cap. 12} --- \emph{Cosa fare}: strategie individuali (livello
  \(\boldsymbol{p}_{\text{personale}}\)), organizzative (livello
  \(\boldsymbol{p}_{\text{azienda/comunità}}\)), statali (livello
  \(\boldsymbol{p}_{\text{nazionale}}\)), condizionate ai quattro mondi
  possibili M1-M4 del Cap. 10.4. La distinzione \emph{protocollo per
  ramo A/B} vs \emph{protocollo per ramo C} è esplicita.
\item
  \textbf{Cap. 13} --- Conclusioni e ricerca futura: posizione del libro
  nella tradizione (Arrighi, Turchin, Wallerstein, Cervo,
  Marcuse-Adorno), agenda di estensione (paper EN, articolo PT, dataset
  50+ paesi), inviti a falsificazione.
\end{itemize}


